\documentclass[../../main.tex]{subfiles}

\begin{document}

\section{Littlewood 的三个原理、Egoroff 定理以及 Lusin 定理}\label{section:ch03-Littlewood-Egoroff-Lusin}

谈到一元实变量的函数论, J. E. Littlewood 说: "所要求的知识范围不像有时料想的那么多. 有三个原理, 大致可用术语表述如下: 每个可测集接近于区间的有限并, 每个可测函数接近于连续函数; 每个逐点收敛的可测函数序列接近于一致收敛的函数序列. 实变函数论中的大部分结果是这些思想相当直观的应用. 学生们掌握了这些等于掌握了大多数情况下实变函数理论所要求的 (知识). 若其中一个原理是解决一个十分真实的问题的明显的方法, 那么自然要问这个 '接近' 是否足够, 而实际上对于一个可以解决的问题, 这个 '接近' 一般是足够的."

第 2 章的\cref{theorem:ch02-有限区间逼近}是 Littlewood 第一原理的精确阐述: 它告诉我们给定有限测度的可测集 $E$, 则对每个 $\varepsilon>0$, 存在开区间的有限不交族, 其并集 $\mathcal{U}$ 在 $m(E\sim\mathcal{U})+m(\mathcal{U}\sim E)<\varepsilon$ 的意义下 "接近等于" $E$.

Littlewood 的最后原理的精确实现是以下令人惊讶的定理.

\begin{theorem}[Egoroff 定理]\label{theorem:ch03-egoroff}
假设 $E$ 具有有限测度. 令 $\{f_{n}\}$ 为 $E$ 上的逐点收敛于实值函数 $f$ 的可测函数序列. 则对每个 $\varepsilon>0$, 存在包含于 $E$ 的闭集 $F$, 使得在 $F$ 上
\begin{align*}
\text{一致地}\{f_{n}\}\to f\text{ 且 }m(E\sim F)<\varepsilon.
\end{align*}
\end{theorem}

为证 Egoroff 定理, 首先建立以下引理是方便的.

\begin{lemma}\label{lemma:ch03-egoroff预备}
在 Egoroff 定理的假设下, 对每个 $\eta>0$ 和 $\delta>0$, 存在 $E$ 的可测子集 $A$ 和指标 $N$, 使得对所有 $n\geqslant N$ 和 $m(E\sim A)<\delta$, 在 $A$ 上 $|f_{n}-f|<\eta$.
\end{lemma}

\begin{proof}
对每个 $k$, 函数 $|f-f_{k}|$ 是恰当定义的, 这是由于 $f$ 是实值的, 它是可测的, 所以集合 $\{x\in E\mid|f(x)-f_{k}(x)|<\eta\}$ 是可测的. 可测集的可数族的交是可测的. 因此
\begin{align*}
E_{n}=\{\text{对所有 }k\geqslant n,\ x\in E\mid|f(x)-f_{k}(x)|<\eta\}
\end{align*}
是一个可测集. 则 $\{E_{n}\}_{n=1}^{\infty}$ 是上升的可测集族, 且 $E=\bigcup_{n=1}^{\infty}E_{n}$, 这是由于 $\{f_{n}\}$ 在 $E$ 上逐点收敛于 $f$. 我们从测度的连续性推出
\begin{align*}
m(E)=\lim_{n\to\infty}m(E_{n}).
\end{align*}
由于 $m(E)<\infty$, 我们可以选取指标 $N$ 使得 $m(E_{N})>m(E)-\delta$. 定义 $A=E_{N}$ 且观察到, 根据测度的分割性质, $m(E\sim A)=m(E)-m(E_{N})<\delta$.
\end{proof}

\begin{proof}[Egoroff 定理的证明]
对每个自然数 $n$, 令 $A_{n}$ 为 $E$ 的可测子集而 $N(n)$ 是满足前一个引理结论的指标, 其中 $\delta=\varepsilon/2^{n+1}$ 而 $\eta=1/n$, 即
\begin{align}
m(E\sim A_{n})<\varepsilon/2^{n+1}\label{align:ch03-A_n估计}
\end{align}
与
\begin{align}
\text{对所有 }k\geqslant N(n),\quad\text{在 }A_{n}\text{ 上 }|f_{k}-f|<1/n.\label{align:ch03-A_n一致}
\end{align}
定义
\begin{align*}
A=\bigcap_{n=1}^{\infty}A_{n}.
\end{align*}
根据 De Morgan 等式、测度的可数次可加性以及\eqref{align:ch03-A_n估计},
\begin{align*}
m(E\sim A)=m\left(\bigcup_{n=1}^{\infty}[E\sim A_{n}]\right)\leqslant\sum_{n=1}^{\infty}m(E\sim A_{n})<\sum_{n=1}^{\infty}\varepsilon/2^{n+1}=\varepsilon/2.
\end{align*}
我们宣称 $\{f_{n}\}$ 在 $A$ 上一致收敛于 $f$. 事实上, 令 $\varepsilon>0$. 选取指标 $n_{0}$ 使得 $1/n_{0}<\varepsilon$. 则由\eqref{align:ch03-A_n一致},
\begin{align*}
\text{对所有 }k\geqslant N(n_{0}),\quad\text{在 }A_{n_{0}}\text{ 上 }|f_{k}-f|<1/n_{0}.
\end{align*}
然而, $A\subseteq A_{n_{0}}$ 且 $1/n_{0}<\varepsilon$, 因此
\begin{align*}
\text{对于 }k\geqslant N(n_{0}),\quad\text{在 }A\text{ 上 }|f_{k}-f|<\varepsilon.
\end{align*}
因此在 $A$ 上 $\{f_{n}\}$ 一致收敛于 $f$ 且 $m(E\sim A)<\varepsilon/2$.

最后, 根据第 2 章的\cref{theorem:ch02-有限区间逼近}, 我们可以选取包含于 $A$ 的闭子集 $F$, 使得 $m(A\sim F)<\varepsilon/2$. 因此 $m(E\sim F)<\varepsilon$ 且在 $F$ 上一致地 $\{f_{n}\}\to f$.
\end{proof}

若收敛是 a.e. 逐点的且极限函数是 a.e. 有限的, 显然 Egoroff 定理也成立.

我们现在在可测函数是简单的情形下给出 Littlewood 第二原理的精确版本, 接着用这个特殊的情形去证明该原理的一般情形——Lusin 定理.

\begin{proposition}\label{proposition:ch03-简单函数连续逼近}
令 $f$ 为定义在 $E$ 上的简单函数. 则对每个 $\varepsilon>0$, 存在 $\mathbf{R}$ 上的连续函数 $g$ 和一个包含于 $E$ 的闭集 $F$, 使得
\begin{align*}
\text{在 }F\text{ 上},\quad f=g\text{ 且 }m(E\sim F)<\varepsilon.
\end{align*}
\end{proposition}

\begin{proof}
令 $a_{1},a_{2},\cdots,a_{n}$ 为 $f$ 取的有限个不同的值, 并且令它们分别在集合 $E_{1},E_{2},\cdots,E_{n}$ 上被取到. 由于这些 $a_{k}$ 是不同的, 族 $\{E_{k}\}_{k=1}^{n}$ 是不交的. 根据第 2 章的\cref{theorem:ch02-有限区间逼近}, 我们可以选取闭集 $F_{1},F_{2},\cdots,F_{n}$ 使得对每个指标 $k$ ($1\leqslant k\leqslant n$),
\begin{align*}
F_{k}\subseteq E_{k}\text{ 且 }m(E_{k}\sim F_{k})<\varepsilon/n.
\end{align*}
则 $F=\bigcup_{k=1}^{n}F_{k}$ 作为闭集的有限族的并集是闭的. 由于 $\{E_{k}\}_{k=1}^{n}$ 是不交的,
\begin{align*}
m(E\sim F)=m\left(\bigcup_{k=1}^{n}[E_{k}\sim F_{k}]\right)=\sum_{k=1}^{n}m(E_{k}\sim F_{k})<\varepsilon.
\end{align*}
在 $F$ 上定义 $g$ 为在 $F_{k}$ 上取值 $a_{k}$ 的函数, $1\leqslant k\leqslant n$. 由于族 $\{F_{k}\}_{k=1}^{n}$ 是不交的, $g$ 是恰当定义的. 此外, $g$ 在 $F$ 上是连续的, 这是由于对点 $x\in F_{i}$, 存在包含 $x$ 的开区间, 它与闭集 $\bigcup_{k\neq i}F_{k}$ 不交, 因而函数 $g$ 在该区间与 $F$ 的交集上是常数. 但 $g$ 可从闭集 $F$ 上的连续函数延拓为整个 $\mathbf{R}$ 上的连续函数（见习题 25）. $\mathbf{R}$ 上的连续函数 $g$ 有我们所要的逼近性质.
\end{proof}

\begin{theorem}[Lusin 定理]\label{theorem:ch03-lusin}
令 $f$ 为定义在 $E$ 上的实值可测函数. 则对每个 $\varepsilon>0$, 存在 $\mathbf{R}$ 上的连续函数 $g$ 和一个包含于 $E$ 的闭集 $F$, 使得
\begin{align*}
\text{在 }F\text{ 上 }f=g\text{ 且 }m(E\sim F)<\varepsilon.
\end{align*}
\end{theorem}

\begin{proof}
我们考虑 $m(E)<\infty$ 的情形, 而将推广到 $m(E)=\infty$ 作为练习. 根据简单逼近定理, 存在定义在 $E$ 上的简单函数序列 $\{f_{n}\}$ 在 $E$ 上逐点收敛于 $f$. 令 $n$ 为自然数. 根据前一个命题, 其中用 $f_{n}$ 代替 $f$ 且用 $\varepsilon/2^{n+1}$ 代替 $\varepsilon$, 我们可以选取 $\mathbf{R}$ 上的连续函数 $g_{n}$ 和包含于 $E$ 的闭集 $F_{n}$, 使得
\begin{align*}
\text{在 }F_{n}\text{ 上 }f_{n}=g_{n}\text{ 且 }m(E\sim F_{n})<\varepsilon/2^{n+1}.
\end{align*}
根据 Egoroff 定理, 存在包含于 $E$ 的闭集 $F_{0}$ 使得 $\{f_{n}\}$ 在 $F_{0}$ 上一致收敛到 $f$ 且 $m(E\sim F_{0})<\varepsilon/2$. 定义 $F=\bigcap_{n=0}^{\infty}F_{n}$. 根据 De Morgan 等式和测度的可数次可加性, 观察到
\begin{align*}
m(E\sim F)=m\left([E\sim F_{0}]\cup\bigcup_{n=1}^{\infty}[E\sim F_{n}]\right)\leqslant\varepsilon/2+\sum_{n=1}^{\infty}\varepsilon/2^{n+1}=\varepsilon.
\end{align*}
由于集合 $F$ 是闭集的交, 它是闭的. 由于 $F\subseteq F_{n}$ 且在 $F_{n}$ 上 $f_{n}=g_{n}$, 每个 $f_{n}$ 在 $F$ 上是连续的. 最后, 由于 $F\subseteq F_{0}$, $\{f_{n}\}$ 在 $F$ 上一致收敛于 $f$. 然而, 连续函数的一致极限是连续的, 因此 $f$ 在 $F$ 上的限制是连续的. 最后, 存在定义在整个 $\mathbf{R}$ 上的连续函数 $g$, 它在 $F$ 上的限制等于 $f$（见习题 25）. 函数 $g$ 有我们所要的逼近性质.
\end{proof}

\begin{exercise}\label{exercise:ch03-25}
假定 $f$ 是在实数的闭集 $F$ 上连续的函数. 证明 $f$ 有到整个 $\mathbf{R}$ 的连续延拓. 这是即将介绍的 Tietze 延拓定理的特殊情形.（提示: 将 $\mathbf{R}\sim F$ 表示为开区间的可数不交族的并集, 且定义 $f$ 在这些区间的每个闭包上是线性的.）
\end{exercise}

\begin{exercise}\label{exercise:ch03-26}
对于 Lusin 定理叙述中的函数 $f$ 和集合 $F$, 证明 $f$ 在 $F$ 上的限制是连续函数. 是否一定存在一些点, 当考虑 $f$ 为 $E$ 上的函数时, 它在这些点是连续的?
\end{exercise}

\begin{exercise}\label{exercise:ch03-27}
证明: 若去掉定义域具有有限测度的假设, Egoroff 定理的结论可能不成立.
\end{exercise}

\begin{exercise}\label{exercise:ch03-28}
证明: 若收敛是 a.e. 逐点的而 $f$ 是 a.e. 有限的, Egoroff 定理仍然成立.
\end{exercise}

\begin{exercise}\label{exercise:ch03-29}
证明 Lusin 定理可推广到 $E$ 具有无穷测度的情形.
\end{exercise}

\begin{exercise}\label{exercise:ch03-30}
证明 Lusin 定理可推广到 $f$ 不必是实值, 但是可以为 a.e. 有限的情形.
\end{exercise}

\begin{exercise}\label{exercise:ch03-31}
令 $\{f_{n}\}$ 为 $E$ 上的可测函数序列, 它在 $E$ 上逐点收敛于实值函数 $f$. 证明 $E=\bigcup_{k=1}^{\infty}E_{k}$, 其中对每个指标 $k$, $E_{k}$ 是可测的, 且若 $k>1$, 在每个 $E_{k}$ 上 $\{f_{n}\}$ 一致收敛于 $f$, 且 $m(E_{1})=0$.
\end{exercise}

\end{document}